% Original indexing diagram based on t10q1.ipynb, cells 29--30 (zero-based).
% Highlighting indicates array selection only; no model scores are plotted.
\begin{tikzpicture}[font=\small,>=Stealth]
  \fill[noteBlue!10] (0,0) rectangle (0.24,2.1);
  \fill[noteBlue!35] (0,1) rectangle (2.6,1.18);
  \draw (0,0) rectangle (2.6,2.1);
  \draw[noteBlue] (0.24,0) -- (0.24,2.1);
  \node[align=center] at (1.3,2.55) {固定头 $i$\\$197\times197$};
  \node[anchor=east,font=\scriptsize] at (0,1.09) {行 100};
  \node[anchor=north,font=\scriptsize,align=center] at (1.3,-0.12)
    {列 0：class token\\列 1--196：patches};
  \draw[->,thick,noteBlue] (2.8,1.09) -- (4.15,1.09)
    node[midway,above,font=\scriptsize] {取 \texttt{1:}};
  \filldraw[fill=noteBlue!25,draw=noteBlue] (4.3,0.99) rectangle (6.9,1.19);
  \node[align=center] at (5.6,1.72) {196 个分数};
  \node[align=center,font=\scriptsize] at (5.6,0.45) {保持原来的 patch 顺序};
  \draw[->,thick,noteBlue] (7.05,1.09) -- (8.45,1.09)
    node[midway,above,font=\scriptsize] {reshape};
  \begin{scope}[shift={(8.65,0.04)}]
    \fill[noteBlue!6] (0,0) rectangle (2.1,2.1);
    \draw[step=0.15,gray!60,thin] (0,0) grid (2.1,2.1);
    \draw (0,0) rectangle (2.1,2.1);
  \end{scope}
  \node[align=center] at (9.7,2.55) {$14\times14$\\patch 位置图};
  \node[align=center,font=\scriptsize] at (9.7,-0.35) {每格对应一个 key patch};
\end{tikzpicture}
